\scanpage{27}

\uline{Remarque} Soient

\diagram{\begin{tikzcd}[ampersand replacement=\&] k \arrow[r,dash] \arrow[d,dash] \& K \arrow[d,dash] \\ k'=k\cap K^p \arrow[r,dash] \& K^p \end{tikzcd}}

alors deux éléments $\deleted{\unsure{z|x}}x,y$ de $k$ peuvent être p-indépendants \inserted{ds $K$ sur} \deleted{\unsure{pour $k$|sur $k$|parfait}} $k\cap K^p$, et non \inserted{dans $K$} sur $K^p$, i.e. $d_{k/k'}x$ et $d_{k/k'}y$ peuvent être lin. indépendants, sans que $d_{K\deleted{\unsure{/|1}}/K^p}x$ et $d_{K/K^p}y$ le soient; \uline{à fortiori}, $\Omega^1_{k/k\cap K^p}\otimes_k K\to\Omega^1_K$ peut ne être pas injectif, ou \deleted{\unsure{encore}} ce qui revient au =, il peut arriver que le noyau \inserted{$N_{K/k}$} de $\Omega^1_k\otimes_k K\unsure{-|\to}\Omega_K$ ne soit pas engendré sur $K$ par les $dx\otimes1$ \inserted{qu'il contient, (i.e. ceux où} $x\in k\cap K^p$).

\margin{$\Omega^1_{k\cap K^p/\unsure{-|\vphantom{k}}}\otimes_{k\cap K^p}K\to\Omega^1_k\otimes_k\deleted{\unsure{k|K}}K\to\Omega^1_K$}

% slanted
\margin{\uline{N.B.} \deleted{\unsure{Alors|Or}} l'existence de $x,y$ comme ci dessus \deleted{\unsure{n'est pas|n'est point}} signifie que $N_{K/k}$ ne \deleted{peut \unsure{être}} \unsure{être|peut être} ``défini sur $k$''! On peut prendre un exemple où il est de dim \uline{1}}

\unsure{et|or} Voici un exemple. Soient $\mathbb{F}$ le corps premier [ou un corps parfait quelconque], $\underline{\underline{x,y,a_i}}\inserted{(1\leq i\leq p-1)}$ \deleted{\unsure{des|deux}} indéterminées sur $\mathbb{F}$, on pose \deleted{$\beta=y-\alpha x$}

\[ \deleted{K^p=\mathbb{F}(x^p,y^p,\unsure{\beta|\beta/\alpha|\alpha})}\qquad\deleted{\text{d'où }K=\unsure{F|\mathbb F}(x,y^{\unsure{1/p|p|\vphantom{p}}},\beta^{1/p})} \]

\[ \deleted{K=}\quad \left\{\begin{aligned} a_0&=y-\sum_{1}^{p-1}a_ix^i\\ a&=x^p\end{aligned}\right. \qquad \text{d'où}\quad y=\sum_{0}^{p-1}\deleted{\unsure{x|a}}\unsure{a_i|a_i^p}x^i \]

\margin{N.B $a$ et les $a_i{}_{1\leq i\leq p-1}$ sont alg. ind. sur $\mathbb F$ \unsure{Alors|dans} \unsure{on|et}}

\[ \begin{aligned}
K^p&=\deleted{\unsure{K|F}}\mathbb{F}(a,a_0,\deleted{\unsure{\beta|a}},\ldots a_{p-1})=\deleted{\unsure{K|\mathbb F}(\alpha,x^p,y-\alpha x)}\\
&\inserted{\mathbb{F}\bigl(x^p,a_1,\ldots a_{p-1},y-\sum_1^{p-1}a_i x^i\bigr)}\\
&\deleted{K=\mathbb{F}(\alpha^{1/p},x,(y-\alpha x)^{1/p})}\quad \text{i.e. } K=\mathbb{F}(x,a_1^{1/p},\ldots a_{p-1}^{1/p},(\ )^{1/p})\\
& k=\mathbb F(x,y)
\end{aligned} \]

\margin{[N.B. $x,\unsure{z|y}\in K$, \unsure{donc|d'où|et} $y=\alpha x+\beta\in K$, donc $k\subset K$}
% archive sheet number: 26
